\Article{18}
\textsc{Problem.} \textit{Let \(U\) denote a given function of the unknown quantities \(V,V',V''\) etc.; the mean error \(M\) to be feared in determining the value of \(U\) is required, if for \(V,V',V''\) etc. we adopt their observed values rather than their true values, the observations being mutually independent and subject respectively to the mean errors \(m,m',m''\) etc.}

\textit{Solution.} Denoting the errors in the observed values of \(V,V',V''\) etc. by \(e,e',e''\) etc., the error thereby arising in the value of \(U\) can be expressed by the linear function
\[
\lambda e+\lambda'e'+\lambda''e''+\text{etc.}=E
\]
where \(\lambda,\lambda',\lambda''\) etc. are the values of the differential quotients \(\frac{\mathrm{d}U}{\mathrm{d}V},\frac{\mathrm{d}U}{\mathrm{d}V'},\frac{\mathrm{d}U}{\mathrm{d}V''}\) etc. for the true values of \(V,V',V''\) etc., provided the observations are sufficiently exact to allow the squares and products of the errors to be neglected. From this it follows first, since the errors of the observations are supposed free of constant parts, that the mean value of \(E\) is \(=0\). Further, the mean error to be feared in the value of \(U\) will be the square root of the mean value of \(EE\), or \(MM\) will be the mean value of the sum
\[
\begin{aligned}
\lambda\lambda ee+\lambda'\lambda'e'e'+\lambda''\lambda''e''e''+\text{etc.}{}\\
{}+2\lambda\lambda'ee'+2\lambda\lambda''ee''+2\lambda'\lambda''e'e''+\text{etc.}
\end{aligned}
\]
But the mean value of \(\lambda\lambda ee\) becomes \(\lambda\lambda mm\), the mean value of \(\lambda'\lambda'e'e'\) becomes \(=\lambda'\lambda'm'm'\) etc.; finally, the mean values of the products \(2\lambda\lambda'ee'\) etc. all become \(=0\). Hence we conclude
\[
M=\BareRoot{(\lambda\lambda mm+\lambda'\lambda'm'm'+\lambda''\lambda''m''m''+\text{etc.})}
\]
\SourceLeaf{025}\par
It will be appropriate to add some remarks to this solution.

I. Insofar as, by regarding the errors of observations as quantities of the first order, quantities of higher orders are neglected, we may also adopt in our formula for \(\lambda,\lambda',\lambda''\) etc. those values of the quotients \(\frac{\mathrm{d}U}{\mathrm{d}V}\) etc. that arise from the observed values of the quantities \(V,V',V''\) etc. Whenever \(U\) is a linear function, there will evidently be no difference at all.

II. If we prefer to introduce the weights of the observations instead of their mean errors, let these, according to an arbitrary unit, be respectively \(p,p',p''\) etc., and let \(P\) be the weight of the determination of the value of \(U\) arising from the observed values of the quantities \(V,V',V''\) etc. We shall thus have
\[
P=\frac{1}{\dfrac{\lambda\lambda}{p}+\dfrac{\lambda'\lambda'}{p'}+\dfrac{\lambda''\lambda''}{p''}+\text{etc.}}
\]

III. If \(T\) is another given function of the quantities \(V,V',V''\) etc., and if, for their true values,
\[
\frac{\mathrm{d}T}{\mathrm{d}V}=\kappa,\quad\frac{\mathrm{d}T}{\mathrm{d}V'}=\kappa',\quad\frac{\mathrm{d}T}{\mathrm{d}V''}=\kappa''\text{ etc.}
\]
the error in the determination of the value of \(T\) obtained from the observed values of \(V,V',V''\) etc. will be \(=\kappa e+\kappa'e'+\kappa''e''+\text{etc.},\ =E'\), and the mean error to be feared in this determination \(=\BareRoot{(\kappa\kappa mm+\kappa'\kappa'm'm'+\kappa''\kappa''m''m''+\text{etc.})}\). But the errors \(E,E'\) will evidently no longer be independent of one another. The mean value at issue is that of the product \(EE'\). Unlike the mean value of the product \(ee'\), it will not be \(=0\), but \(=\kappa\lambda mm+\kappa'\lambda'm'm'+\kappa''\lambda''m''m''+\text{etc.}\)

IV. Our problem can also be extended to the case where the values of the quantities \(V,V',V''\) etc. are not found directly by observations but are derived in any way from combinations of observations, provided only that their individual determinations are mutually independent, i.e. based on different observations: whenever this condition does not hold, the formula for \(M\) would become erroneous. For example, if one or another observation that served to determine the value of \(V\) had also been used to determine the value of \(V'\), the errors \(e\) and \(e'\) would no longer be mutually independent, nor therefore would the mean value of the product \(ee'\) be \(=0\). But if in such a case the connection of the quantities \(V,V'\) with the individual observations from which they have been derived is duly considered, the \SourceLeaf{026}mean value of the product \(ee'\) can be assigned with the aid of remark III., and the formula for \(M\) can thus be made complete.

\Article{19}
Let \(V,V',V''\) etc. be functions of the unknowns \(x,y,z\) etc., the number of the former \(=\pi\), the number of the unknowns \(=\rho\); and let us suppose that the values of the functions have been found by observations, either directly or indirectly, as \(V=L,V'=L',V''=L''\) etc., yet in such a way that these determinations were mutually independent. If \(\rho\) is greater than \(\pi\), finding the unknowns evidently becomes an undetermined problem; the next case is where \(\rho\) is equal to \(\pi\). Each of \(x,y,z\) etc. can be reduced, or conceived as reduced, to the form of functions of \(V,V',V''\) etc., so that the values of the former can be found from the observed values of the latter, and at the same time, with the aid of the preceding article, the relative precision to be assigned to each of these determinations can be specified; finally, if \(\rho\) is less than \(\pi\), each of \(x,y,z\) etc. can be reduced in infinitely many different ways to the form of functions of \(V,V',V''\) etc., and thus their values can be obtained in infinitely many different ways. These determinations should indeed agree exactly if the observations had absolute precision; since the case is otherwise, different methods will give different values, and likewise the determinations obtained from different combinations will have unequal precision.

Moreover, if in the second or third case the functions \(V,V',V''\) etc. were such that \(\pi-\rho+1\) of them, or more, could be regarded as functions of the remaining ones, the problem would still be more than determined with regard to the latter functions, but undetermined with regard to the unknowns \(x,y,z\) etc.; their values could not be determined even if absolutely exact values of the functions \(V,V',V''\) etc. were given: but we shall exclude this case from our investigation.

Whenever \(V,V',V''\) etc. are not themselves \textit{linear} functions of their indeterminates, this will be achieved if, in place of the original unknowns, we introduce their differences from approximate values. These may be supposed known from elsewhere. We shall denote the mean errors to be feared in the determinations \(V=L,V'=L',V''=L''\) etc. respectively by \(m,m',m''\) etc., and the weights of the determinations by \(p,p',p''\) etc., so that \(pmm=p'm'm'=p''m''m''\) etc. We shall suppose the ratio of the mean errors to one another known, so that the weights, one of which can be chosen at will, are known. Finally, we shall put
\SourceLeaf{027}
\[
\begin{aligned}
(V-L)\BareRoot{p}&=v,\quad(V'-L')\BareRoot{p'}=v',\\
(V''-L'')\BareRoot{p''}&=v''\text{ etc.}
\end{aligned}
\]
Evidently, therefore, the case will be the same as if direct observations of equal precision, namely whose mean error is \(=m\BareRoot{p}=m'\BareRoot{p'}=m''\BareRoot{p''}\) etc., or to which the weight \(=1\) is assigned, had given
\[
v=0,\quad v'=0,\quad v''=0\text{ etc.}
\]

\clearpage
